\section{Definitions}\label{sec:tsps-definitions}

Structure-preserving signatures sign vectors of messages that lie in
pairing-friendly groups, rather than arbitrary strings, and the corresponding
verification algorithm checks equations comprising products of pairings.

\subsection{Structure-preserving signatures}

A structure-preserving signature scheme is defined by the following
algorithms:

\begin{description}

  \item[$(\pk, \sk) \sample \SPSgen(\secparam, \ell)$ :] The key generation
    algorithm, on input security parameter $\secparam$ and vector length $\ell$,
    outputs key pair $(\pk, \sk)$. Public keys are group elements.

  \item[$\signature \sample \SPSsign(\sk, \vec{{m}})$ :] The signing algorithm, on
    input secret key $\sk$ and message $\vec{{m}}$, outputs
    signature $\signature$. Messages and signatures are vectors of group elements, 
    messages from the space $\group{1}^\ell$, and signatures in 
    $\group{1}^{poly(\secparam)} \times \group{2}^{poly(\secparam)}$.

  \item[$\{0,1\} \leftarrow \SPSverify(\pk, \vec{{m}}, \signature)$ :] The
    verification algorithm, on input public key $\pk$, message $\vec{{m}}$,
    and signature $\signature$, returns 0 or 1. Verification consists only of 
    evaluating pairing product equations.

\end{description}

A structure-preserving signature scheme satisfies correctness if for all
security parameters $\secparam$, all $(\pk, \sk) \sample \SPSgen(\secparam,
l)$, and all $\vec{{m}} \in \group{1}^\ell$, we have $\SPSverify(\pk, \vec{{m}},
\SPSsign(\sk, \vec{{m}})) = 1$.

\paragraph{Unforgeability.} For a PPT adversary $\adv$, define advantage as
follows:
%
\begin{equation*}
%
  \Adv_{\mathrm{SPS}, \adv}^{\text{UF-CMA}}(\secparam) =
  \Pr\left\lbrack \mathbf{G}_{\mathrm{SPS}, \adv}^{\text{UF-CMA}}(\secparam) =
  1\right\rbrack .
%
\end{equation*}
%
Here, game $\mathbf{G}_{\mathrm{SPS},\adv}^{\text{UF-CMA}}$ and signing oracle
$\signo$ are given in Experiment~\ref{exp:sps-uf-cma}. A structure-preserving
signature scheme achieves UF-CMA security if, for every PPT adversary $\adv$,
the following holds:
%
\begin{equation*}
%
  \Adv_{\mathrm{SPS}, \adv}^{\text{UF-CMA}}(\secparam) \leq
  \operatorname{negl}(\secparam).
%
\end{equation*}
%
Excluding the message $\vec{{m}^*}$ from the set $S$ of
queried messages, as shown in Experiment~\ref{exp:sps-uf-cma}, gives
weak unforgeability (wUF-CMA), the notion we rely on throughout this
chapter. Excluding instead the pair $(\vec{{m}^*}, \signature^*)$ gives 
strong unforgeability (sUF-CMA), which
rules out an adversary producing a fresh signature on an already-signed
message.

\begin{experimentbox}{Unforgeability under chosen message attack for
  structure-preserving signatures}{sps-uf-cma}

  \begin{enumerate}

    \item The challenger generates $(\pk, \sk) \gets \SPSgen(\secparam,
      l)$ and sends $\pk$ to $\adv$.

    \item The adversary is given oracle access to $\signo$, which on input
      $\vec{{m}}$, computes $\signature \gets \SPSsign(\sk, \vec{{m}})$, and if
      $\signature \neq \perp$, adds $\vec{{m}}$ to $S$, before returning $\signature$.

    \item The adversary outputs $(\vec{{m}^*}, \signature^*)$.

  \end{enumerate}

  $\adv$ wins the experiment if $\SPSverify(\pk, \vec{{m}^*}, \signature^*) =
  1$ and $\vec{{m}^*} \notin S$.

\end{experimentbox}

\subsection{Threshold structure-preserving signatures}

A threshold structure-preserving signature scheme is a structure-preserving
signature scheme whose signing key is distributed across multiple parties. It
is defined by algorithms $\TSPSgen$, $\TSPSsign$, $\TSPScombine$, and
$\TSPSverify$ as follows.

\begin{description}

  \item[$(\pk, \lbrack \sk\rbrack ) \sample \TSPSgen(l, t, n)$ :] The key
    generation algorithm takes a triple of integers $(l, t, n > t)$, where
    $l$ is the length of the vectors to be signed, $t$ is an upper bound on
    the number of corrupted parties, and $n$ is the total number of
    threshold signers. It returns a public key $\pk$ and $n$ shares
    $\lbrack \sk\rbrack $ such that any $t+1$ shares reconstruct the
    secret key underlying $\pk$; $\pk$ consists of group elements, and $\sk$
    of elements of $\Zp$ and/or group elements.

  \item[$\lbrack \signature\rbrack  \sample \TSPSsign(\lbrack \sk\rbrack , \vec{{m}})$ :] The threshold signing algorithm takes a
    secret key share $\lbrack \sk\rbrack $ and a vector of messages
    $\vec{{m}}$ with elements in $\group{1}$ and/or $\group{2}$, and
    outputs a signature share $\lbrack \signature\rbrack $ with elements in
    $\group{1}$ and $\group{2}$.

  \item[$\signature \leftarrow \TSPScombine(\lbrack \signature\rbrack )$ :] The
    combining algorithm is non-interactive, and takes signature shares
    $\lbrack \signature\rbrack $ to output signature $\signature$.

  \item[$\{0,1\} \leftarrow \TSPSverify(\pk, \vec{{m}}, \signature)$ :] The
    verification algorithm takes public key $\pk$, vector of messages
    $\vec{{m}}$, and signature $\signature$, and returns 0 or 1. $\TSPSverify$ 
    is equal to $\SPSverify$.

\end{description}

\paragraph{Unforgeability.} As with structure-preserving signatures, we
define the advantage of a PPT adversary $\adv$, though threshold
structure-preserving signatures admit a hierarchy of increasingly strong
notions of unforgeability, distinguishing fully and partially
non-interactive schemes~\cite{C:BCKMTZ22}. With TS-UF-0 security, the
adversary has access to a partial signature oracle $\signo$, and the
unforgeability game is defined so that the adversary cannot win with
message-signature pair $(\vec{{m}^*}, \signature^*)$ if the partial
signature oracle has already been queried, even once, on $\vec{{m}^*}$.
TS-UF-1 relaxes this: the winning condition can be triggered as long as the
adversary has queried for partial signatures on $\vec{{m}^*}$ from no more
than $t - |\Cor| - 1$ parties, where $\Cor$ is the set of corrupted parties.
When the adversary corrupts $t-1$ parties, TS-UF-0 and TS-UF-1 coincide. In
the adaptive variants, adp-TS-UF-0 and adp-TS-UF-1, the adversary
additionally has access to a corruption oracle $\corrupto$, alongside the
partial signature oracle $\signo$. We rely on adaptive TS-UF-1 security
throughout this chapter.

Let $\mathrm{TSPS} = (\algo{Setup}, \algo{KeyGen}, \algo{OffSign}, \algo{OnSign}, \algo{Combine},
\SPSverify)$ be an $(n, t)$-TSPS scheme over message space $\mathcal{M}$. The
advantage of a PPT adversary $\adv$ is defined as follows:
%
\begin{equation*}
%
  \Adv_{\mathrm{TSPS}, \adv}^{\text{adp-TS-UF-1}}(\secparam) =
  \Pr\left\lbrack \mathbf{G}_{\mathrm{TSPS}, \adv}^{\text{adp-TS-UF-1}}(\secparam) = 1
  \right\rbrack .
%
\end{equation*}
%
Here, game $\mathbf{G}_{\mathrm{TSPS}, \adv}^{\text{adp-TS-UF-1}}$ and its
oracles $\signo$ and $\corrupto$ are given in
Experiment~\ref{exp:tsps-adp-ts-uf-1}. A threshold structure-preserving
signature scheme achieves adaptive TS-UF-1 security if, for every PPT
adversary $\adv$, the following holds:
%
\begin{equation*}
%
  \Adv_{\mathrm{TSPS}, \adv}^{\text{adp-TS-UF-1}}(\secparam)
  \leq \operatorname{negl}(\secparam).
%
\end{equation*}

\begin{experimentbox}{Adaptive unforgeability under chosen message attack for
  threshold structure-preserving signatures (adp-TS-UF-1)}{tsps-adp-ts-uf-1}

  \begin{enumerate}

    \item The challenger generates $\algo{pp} \gets \algo{Setup}(1^\secparam)$
      and sends it to $\adv$.

    \item The adversary outputs $(n, t, \Cor)$, and the challenger sets
      $\mathcal{H} := [1,n] \backslash \Cor$.

    \item The challenger generates $(\pk, \lbrack \sk\rbrack , \lbrack \pk\rbrack ) \gets \algo{KeyGen}(l, t, n)$, and sends
      $\pk$, $\{\sk_i\}_{i \in \Cor}$, and $\lbrack \pk\rbrack $ to
      $\adv$.

    \item The adversary is given oracle access to $\signo$, which on input
      $(i, \vec{{m}})$, asserts $\vec{{m}} \in \mathcal{M} \wedge i \in
      \mathcal{H}$, computes $\sigma_i \gets \algo{ThreshSign}(\sk_i,
      \vec{{m}})$, and if $\sigma_i \neq \perp$, adds $i$ to $S(\vec{{m}})$,
      before returning $\sigma_i$.

    \item The adversary is given oracle access to $\corrupto$, which on
      input $k$, returns $\perp$ if $k \in \Cor$, and otherwise adds $k$ to
      $\Cor$, removes $k$ from $\mathcal{H}$, and returns $\sk_k$.

    \item The adversary outputs $(\vec{{m}^*}, \signature^*)$.

  \end{enumerate}

  $\adv$ wins the experiment if $\SPSverify(\pk, \vec{{m}^*}, \signature^*) =
  1$, $|\Cor| < t$, and $|S(\vec{{m}^*})| < t - |\Cor|$, where $S(\vec{{m}})$
  records which honest signers have issued a partial signature on
  $\vec{{m}}$.

\end{experimentbox}

Let $x \in \Zp$ be a secret value additively shared between $n$ parties such
that any $t+1$ parties can reconstruct it with the help of a pre-defined
linear combination. We denote such a sharing $\lbrack x\rbrack $: the
vector $(x_1, \ldots, x_n)$ where $x_i$ is the share held by the
$i^{\mathrm{th}}$ party. For $\vec{x} = (x_1, \ldots, x_l) \in \Zp^l$, we
write $\lbrack \vec{x}\rbrack $ for the vector $(\lbrack x_1\rbrack , \ldots, \lbrack x_l\rbrack )$.

Our threshold signing algorithms are built from secure multiparty
computation (MPC). We introduce the ideal functionalities for MPC that they
rely on below.

\subsection{Secure multiparty computation}\label{sec:MPC}

Our threshold signatures are built in a hybrid model: rather than specifying
how the signers jointly compute each step, we assume access to ideal
functionalities that perform the relevant computation on their behalf, and
describe our schemes in terms of calls to these functionalities. This
section describes the ideal functionalities we assume for secure MPC over a
prime order field $\Zp$ and a generic prime order group $\G$.
Section~\ref{sec:tsps-evaluation} then summarises concrete instantiations of
these ideal functionalities in the honest- and dishonest-majority settings,
via ATLAS and MASCOT respectively.

\subsubsection{Functionalities for secure multiparty computation in $\Zp$}

\begin{figure*} \centering \hspace*{-0.05\textwidth}
  \begin{minipage}{1.1\textwidth}

    \begin{framed} \scriptsize

      \setlist[enumerate]{ topsep=2pt, partopsep=0pt, parsep=0pt, itemsep=1pt }

      \begin{multicols}{2}

        {$\idfu{\algo{RandShare}}$.} On receiving a fresh identifier
        $\mathit{id}_x$ from the honest parties, contact the adversary with
        this identifier. In response, receive the set of shares
        $\{x^*_i\}_{{P}_i \in \Cor}$, then:

        \begin{enumerate}

          \item Sample a random $t$-additive sharing $\lbrack x\rbrack 
            = (x_1, ..., x_n)$ such that $x^*_i = x_i, \forall i : {P}_i \in
            \Cor$. If no such sharing exists, send $\bot$ to all honest
            parties and stop.

          \item Otherwise, distribute to each honest party its share
            $\langle \mathit{id}_x, x_i\rangle$, and store entries $\langle
            \algo{share}, \mathit{id}_x, x_i \rangle$ and $\langle
            \algo{secret}, \mathit{id}_x, x \rangle$.

        \end{enumerate}

        {$\idfu{\algo{Add}}$.} On receiving from the honest parties tuples
        $(\mathit{id}_z, \langle \mathit{id}_x, x_i\rangle, \langle
        \mathit{id}_y,y_i\rangle)_{{P}_i \in \mathcal{H}}$, do:

        \begin{enumerate}

          \item Retrieve all entries $\langle \algo{share}, \mathit{id}_x,
            x_i\rangle$ and $\langle \algo{share}, \mathit{id}_y,
            y_i\rangle$, followed by entries $\langle \algo{secret},
            \mathit{id}_x, x\rangle$ and $\langle \algo{secret},
            \mathit{id}_y, y\rangle$, compute $z_i = x_i+y_i$ and $z = x+y$,
            and store entries $\langle \algo{share}, \mathit{id}_z,
            z_i\rangle$ and $\langle \algo{secret}, \mathit{id}_z, z
            \rangle$.

          \item Output to each honest party ${P}_i$ its share $z_i$.

        \end{enumerate}

        {$\idfu{\algo{Mul}}$.} On receiving from the honest parties tuples
        $(\mathit{id}_z, \langle \mathit{id}_x, x_i\rangle, \langle
        \mathit{id}_y,y_i\rangle)_{{P}_i \in \mathcal{H}}$ and from the
        adversary tuples $(\mathit{id}_z, \langle \mathit{id}_x, x_i
        \rangle, \langle \mathit{id}_y, y_i\rangle)_{{P}_i \in \Cor}$, do:

        \begin{enumerate}

          \item Retrieve entries $\langle \algo{share}, \mathit{id}_x,
            x_i\rangle$ and $\langle \algo{share}, \mathit{id}_y, y_i
            \rangle$ using the values $x_i$ and $y_i$ provided in the
            previous step.

          \item If at least $t+1$ entries are successfully retrieved for
            $\mathit{id}_x$ and also for $\mathit{id}_y$, then:

            \begin{enumerate}
              \item fetch entries $\langle \algo{secret}, \mathit{id}_x,x\rangle$
                and $\langle \algo{secret}, \mathit{id}_y,y\rangle$;
              \item compute $z = xy$ and a $t$-additive sharing $\lbrack z
\rbrack  = (z_1, ..., z_n)$ of $z$;
              \item send to each party ${P}_i$ its share $z_i$;
              \item and store entries $\langle \algo{share}, \mathit{id}_z,
                z_i\rangle$, $\forall i \in [n]$ and entry $\langle
                \algo{secret}, \mathit{id}_z, z\rangle$.
            \end{enumerate}

          \item Otherwise, send $\bot$ to the honest parties and stop.

        \end{enumerate}

        {$\idfu{\algo{Inv}}$.} On receiving from the honest parties their
        tuples $(\mathit{id}_z, \langle \mathit{id}_x,
        x_i\rangle)_{{P}_i \in \mathcal{H}}$ and from the adversary the
        tuples $(\mathit{id}_z, \langle \mathit{id}_x, x_i
        \rangle)_{{P}_i \in \Cor}$, do:

        \begin{enumerate}

          \item Retrieve entries $\langle \algo{share}, \mathit{id}_x,
            x_i\rangle$ using the values $x_i$ received in the previous
            step.

          \item If at least $t+1$ entries are successfully retrieved, then:

            \begin{enumerate}
              \item get entry $\langle \algo{secret}, \mathit{id}_x, x
                \rangle$;
              \item compute $z = 1/x$ and a $t$-additive sharing
                $\lbrack z\rbrack  =(z_1, ..., z_n)$ of $z$;
              \item send to each party ${P}_i$ its share $z_i$;
              \item and store entries $\langle \algo{share}, \mathit{id}_z,
                z_i\rangle$, $\forall \ i \in [n]$ and entry $\langle
                \algo{secret}, \mathit{id}_z, z\rangle$.
            \end{enumerate}

          \item Otherwise, send $\bot$ to the honest parties and stop.

        \end{enumerate}

        {$\idfu{\algo{Combine}}$.} On receiving from the honest parties
        their shares $(\langle \mathit{id}_x,x_i\rangle)_{{P}_i \in
        \mathcal{H}}$ and from the adversary the shares $(\langle
        \mathit{id}_x, x_i \rangle)_{{P}_i \in \Cor}$, do:

        \begin{enumerate}

          \item Retrieve entries $\langle \algo{share}, \mathit{id}_x,
            x_i\rangle$ using the values $x_i$ received in the previous
            step.

          \item If at least $t+1$ entries are successfully retrieved, get
            entry $\langle \algo{secret}, \mathit{id}_x, x\rangle$ and send
            $x$ to the honest parties and the adversary.

          \item Otherwise, send $\bot$ to the honest parties and stop.

        \end{enumerate}

      \end{multicols}

    \end{framed}

  \end{minipage}
\caption{Ideal functionalities for MPC in $\Zp$.}
\label{fig:ideal-func}
\end{figure*}

Our constructions of threshold structure-preserving signatures build upon
existing frameworks for secure MPC over inputs in $\Zp$. Concretely, we assume
some realisation of the ideal MPC functionality over $\Zp$ described in
Figure~\ref{fig:ideal-func}, where $\mathcal{H}$ and $\Cor$ denote the sets of honest
and corrupted parties respectively. Unless otherwise specified, we assume
security against malicious corruption of parties. Recall that $\lbrack x
\rbrack $, for $x \in \Zp$, denotes that $x$ is secret-shared: no individual
party has access to $x$, but each party holds some share $x_i$ of $x$ (for
simplicity, we assume this notation incorporates the additional
authentication components required for malicious security).

\paragraph{Linearity-preservation.} Fundamentally, we require that the
secret-shared representation $\lbrack x\rbrack $ is
linearity-preserving: for any $x,y,z,\alpha,\beta\in \Zp$ such that $u =
\alpha\cdot x + \beta \cdot y + z$, given secret shares $\lbrack x\rbrack $ and $\lbrack y\rbrack $
and public values $z,\alpha,\beta$,
the parties can compute a secret sharing of $u$ ``for free'' as follows:
\[\lbrack u\rbrack = \alpha\cdot \lbrack x\rbrack + \beta \cdot
\lbrack y\rbrack + z.\] Since we assume malicious security, this
property must also be preserved for any additional authentication
components, which we otherwise abstract out for ease of exposition.

\paragraph{Functionalities for secure MPC in $\Zp$.} Figure~\ref{fig:ideal-func}
formally describes the ideal functionalities we assume for secure MPC in
$\Zp$. We assume a randomised functionality that generates a secret-shared
representation $\lbrack r\rbrack $ for a randomly sampled value $r\gets
\Zp$, together with four deterministic functionalities, described below
(all arithmetic is in $\Zp$):

\begin{enumerate}

	\item A functionality that adds secret-shared inputs: given $\lbrack x
\rbrack $ and $\lbrack y\rbrack $, produces $\lbrack z
\rbrack $ such that $z = x + y$. By linearity-preservation, this comes
    for free, as it can be performed locally by each party.

	\item A functionality that multiplies secret-shared inputs: given
    $\lbrack x\rbrack $ and $\lbrack y\rbrack $, produces
    $\lbrack z\rbrack $ such that $z = x \cdot y$.

	\item A functionality that inverts a secret-shared value: given
    $\lbrack x\rbrack $, produces $\lbrack z\rbrack $ such that $z
    = x^{-1}$.

	\item A functionality that combines the shares of a secret-shared value
    $\lbrack x\rbrack $: reconstructs and distributes $x$ to all, or a
    subset, of the parties.

\end{enumerate}

\subsubsection{Functionalities for secure multiparty computation in $\G$}

\begin{figure*} \centering \hspace*{-0.05\textwidth}
  \begin{minipage}{1.1\textwidth}

    \begin{framed} \scriptsize

      \setlist[enumerate]{ topsep=2pt, partopsep=0pt, parsep=0pt, itemsep=1pt }

      \begin{multicols}{2}

        {$\idfu{\algo{RandShareGrp}}$.} On receiving a fresh identifier
        $\mathit{id}_{{X}}$ from the honest parties, contact the adversary
        with this identifier. In response, receive the set of shares
        $\{x_i \}_{{P}_i \in \Cor}$, then:

        \begin{enumerate}

          \item Sample a random $t$-additive sharing $\lbrack x\rbrack 
            = (x_1, ..., x_n)$ such that $x_i = x^*_i, \forall i : {P}_i \in
            \Cor$. If no such sharing exists, send $\bot$ to all honest
            parties and stop.

          \item Otherwise:

            \begin{enumerate}
              \item compute the shares $\lbrack {X}\rbrack  = ({X}_1,
                ..., {X}_n)$ such that ${X}_i = x_i\gen$;
              \item distribute to each honest party the tuple $\langle
                \mathit{id}_{{X}}, X_i, x_i\rangle$;
              \item and store entries $\langle \algo{share},
                \mathit{id}_{{X}}, {X}_i \rangle$ and $\langle \algo{secret},
                \mathit{id}_{{X}}, {X} \rangle$.
            \end{enumerate}

        \end{enumerate}

        {$\idfu{\algo{MulGrp}}$.} On receiving from the honest parties
        tuples $(\mathit{id}_{{Z}}, \langle \mathit{id}_{{X}},
        {X}_i\rangle, \langle \mathit{id}_{{Y}},{Y}_i\rangle)_{{P}_i \in
        \mathcal{H}}$, do:

        \begin{enumerate}

          \item Retrieve all entries $\langle \algo{share},
            \mathit{id}_{{X}}, {X}_i\rangle$ and $\langle \algo{share},
            \mathit{id}_{{Y}}, {Y}_i\rangle$, followed by entries $\langle
            \algo{secret}, \mathit{id}_{{X}}, {X}\rangle$ and $\langle
            \algo{secret}, \mathit{id}_{{Y}}, {Y}\rangle$, compute ${Z}_i =
            {X}_i+{Y}_i$ and ${Z} = {X}+{Y}$, and store entries $\langle
            \algo{share}, \mathit{id}_{{Z}}, {Z}_i\rangle$ and $\langle
            \algo{secret}, \mathit{id}_{{Z}}, {Z} \rangle$.

          \item Output to each honest party ${P}_i$ its share ${Z}_i$.

        \end{enumerate}

        {$\idfu{\algo{ExpGenGrp}}$.} On receiving from the honest parties
        tuples $(\mathit{id}_{{Z}}, \langle \mathit{id}_{x},
        x_i\rangle, {Y})_{{P}_i \in \mathcal{H}}$, do:

        \begin{enumerate}

          \item Retrieve all entries $\langle \algo{share}, \mathit{id}_x,
            x_i\rangle$ and entry $\langle \algo{secret}, \mathit{id}_x,
            x\rangle$, compute ${Z}_i = x_i{Y}$ and ${Z} = x{Y}$, and store
            entries $\langle \algo{share}, \mathit{id}_{{Z}}, {Z}_i\rangle$
            and $\langle \algo{secret}, \mathit{id}_{{Z}}, {Z} \rangle$.

          \item Output to each honest party ${P}_i$ its share ${Z}_i$.

        \end{enumerate}

        {$\idfu{\algo{ExpGrp}}$.} On receiving from the honest parties the
        tuples $(\mathit{id}_{{Z}}, \langle \mathit{id}_x, x_i\rangle,
        \langle \mathit{id}_{{Y}},{Y}_i\rangle)_{{P}_i \in \mathcal{H}}$ and
        from the adversary the shares $(\mathit{id}_{{Z}}, \langle
        \mathit{id}_x, x_i \rangle, \langle \mathit{id}_{{Y}},
        {Y}_i\rangle)_{{P}_i \in \Cor}$, do:

        \begin{enumerate}

          \item Retrieve entries $\langle \algo{share}, \mathit{id}_x,
            x_i\rangle$ and $\langle \algo{share}, \mathit{id}_{{Y}}, {Y}_i
            \rangle$ using the values received in the previous step.

          \item If at least $t+1$ entries are successfully retrieved for
            $\mathit{id}_x$ and also for $\mathit{id}_{{Y}}$, then:

            \begin{enumerate}
              \item retrieve entries $\langle \algo{secret},
                \mathit{id}_x,x\rangle$ and $\langle \algo{secret},
                \mathit{id}_{{Y}},{Y}\rangle$;
              \item compute ${Z} = x{Y}$, compute $t$-additive shares
                $(x'_1, ..., x'_n)$ of $x$, and set ${Z}_i = x'_i{Y}$;
              \item send to each party ${P}_i$ its share ${Z}_i$ of ${Z}$;
              \item and store entries $\langle \algo{share},
                \mathit{id}_{{Z}}, {Z}_i\rangle$ and entry $\langle
                \algo{secret}, \mathit{id}_{{Z}}, {Z}\rangle$.
            \end{enumerate}

          \item Otherwise, send $\bot$ to the honest parties and stop.

        \end{enumerate}

        {$\idfu{\algo{CombineGrp}}$.} On receiving from the honest parties
        their shares $(\langle \mathit{id}_{{X}}, {X}_i\rangle)_{{P}_i \in
        \mathcal{H}}$ and from the adversary the shares $(\langle
        \mathit{id}_{{X}}, {X}_i \rangle)_{{P}_i \in \Cor}$, do:

        \begin{enumerate}

          \item Retrieve entries $\langle \algo{share}, \mathit{id}_{{X}},
            {X}_i\rangle$ using the values ${X}_i$ received in the previous
            step.

          \item If at least $t+1$ entries are retrieved successfully, fetch
            entry $\langle \algo{secret}, \mathit{id}_{{X}}, {{X}}\rangle$
            and send ${X}$ to the honest parties and the adversary.

          \item Otherwise, send $\bot$ to the honest parties and stop.

        \end{enumerate}

      \end{multicols}

    \end{framed}

  \end{minipage}
\caption{Ideal functionalities for MPC in $\G$.}
\label{fig:ideal-func-group}
\end{figure*}

We additionally assume a framework for secure MPC over inputs in a cyclic
group $(\G,+)$ of prime order $p$. In practice, this framework is
satisfied by both $\group{1}$ and $\group{2}$, but we present it here in the
abstract to avoid repetition. A group element is shared by sharing a random
exponent $x$ and raising a public generator $\gen$ to each share of $x$; each
resulting group element is then an additive share of $x\gen$, as captured by
the $\idfu{\algo{RandShareGrp}}$ interface below. Concretely, we assume some
realisation of the ideal MPC functionality over $\G$ described in
Figure~\ref{fig:ideal-func-group}, where $\mathcal{H}$ and $\Cor$ again denote the
sets of honest and corrupted parties respectively. As for MPC over $\Zp$,
unless otherwise specified, we assume security against malicious corruption
of parties.

We use the representation $\lbrack X\rbrack $, for $X\in \G$, to
denote that $X$ is secret-shared: no individual party has access to $X$, but
each party holds some share $X_i$ of $X$ (again, we assume this notation
incorporates the authentication components required for malicious security).

\paragraph{Linearity-preservation.} We again require that the secret-shared
representation $\lbrack X\rbrack $ is linearity-preserving: for any
$X,Y,Z\in \G$ and $\alpha,\beta\in \Zp$ such that $U = \alpha X + \beta Y
+ Z$, given secret shares $\lbrack X\rbrack $ and $\lbrack Y
\rbrack $, public group element $Z$, and public values $\alpha,\beta$, the
parties can compute a secret sharing of $U$ ``for free'' as follows:
\[\lbrack U\rbrack = \alpha\lbrack X\rbrack + \beta \lbrack Y
\rbrack + Z.\] Again, since we assume malicious security, this property
must also be preserved for any additional authentication components.

\paragraph{Functionalities for secure MPC in $\G$.} Figure~\ref{fig:ideal-func-group}
describes the functionalities we assume for secure MPC in $\G$, analogous
to those of Figure~\ref{fig:ideal-func} for $\Zp$. We again assume a
randomised functionality that generates a secret-shared representation
$\lbrack R\rbrack $ for a randomly sampled group element $R\gets
\G$, together with four deterministic functionalities, described below
(all group operations below are written additively):

\begin{enumerate}

	\item A functionality that adds secret-shared group elements: given
    $\lbrack X\rbrack $ and $\lbrack Y\rbrack $, produces
    $\lbrack Z\rbrack $ such that $Z = X + Y$. By linearity-preservation
    this comes for free, as it can be performed locally by each party.

	\item A functionality that multiplies a secret-shared field element with a
    public group element: given $\lbrack x\rbrack $ and public $Y$,
    produces $\lbrack Z\rbrack $ such that $Z = xY$. As with addition,
    this comes for free, as it can be executed locally by each party.

	\item A functionality that multiplies a secret-shared field element with a
    secret-shared group element: given $\lbrack x\rbrack $ and
    $\lbrack Y\rbrack $, produces $\lbrack Z\rbrack $ such that $Z
    = xY$.

	\item A functionality that combines the shares of a secret-shared group
    element $\lbrack X\rbrack $: reconstructs and distributes $X$ to
    all, or a subset, of the parties.

\end{enumerate}

Multiplying a secret-shared group element $\lbrack X\rbrack $ by a
public value $y \in \Zp$, to produce $\lbrack Z\rbrack $ such that $Z =
yX$, can also be performed locally, as this is essentially a repeated
application of the local addition functionality above.

In practice, $\G$ corresponds to either $\group{1}$ or $\group{2}$, most
often represented concretely by elliptic curve groups, whose elements are
coordinate pairs satisfying the curve equation, e.g. $(x_0, y_0)$. Sharing a
group element is not achieved by two independent secret sharings of
the coordinates $x_0$ and $y_0$, that is, $(\lbrack x_0\rbrack , \lbrack 
y_0\rbrack )$, as this would require parties to communicate in order to
compute group operations on the shares. Instead, the group element is shared
as a single unit, $\lbrack (x_0,y_0)\rbrack $, so that group operations
can be performed by each party locally, as desired.

